\documentclass[12pt,a4paper]{article}
\usepackage[utf8]{inputenc}
\usepackage[T1]{fontenc}
\usepackage[spanish]{babel}
\usepackage{geometry}
\usepackage{hyperref}
\usepackage{longtable}
\usepackage{xcolor}
\usepackage{listings}
\usepackage{graphicx}
\usepackage{float}

\geometry{margin=2.5cm}
\hypersetup{colorlinks=true, linkcolor=blue, urlcolor=blue}

\lstdefinestyle{terminal}{
  basicstyle=\ttfamily\small,
  breaklines=true,
  frame=single,
  backgroundcolor=\color{gray!8}
}

\newcommand{\evidenceimage}[3]{
  \begin{figure}[H]
    \centering
    \IfFileExists{figures/#1}{\includegraphics[width=0.92\textwidth]{figures/#1}}{\fbox{\parbox[c][5cm][c]{0.86\textwidth}{\centering Captura pendiente: \texttt{figures/#1}}}}
    \caption{#2}
    \label{#3}
  \end{figure}
}

\title{Manual de Despliegue Local\\Sistema de Reserva de Canchas Deportivas}
\author{Enzo Aliatis \and Jordan Murillo \and Galo Suarez \and Mauro Cadme}
\date{Agosto 2026}

\begin{document}
\maketitle

\section{Objetivo}

Este manual describe como levantar localmente el sistema completo de reserva de canchas deportivas. El entorno integra cuatro microservicios Spring Boot, cuatro bases PostgreSQL, un shell Angular y tres microfrontends remotos mediante Module Federation.

\section{Prerequisitos}

\begin{itemize}
  \item Docker y Docker Compose.
  \item Puertos libres: 4300, 4201, 4202, 4203, 8081, 8082, 8083 y 8084.
  \item Git para clonar el repositorio.
\end{itemize}

\section{Servicios desplegados}

\begin{longtable}{p{4cm}p{3cm}p{6cm}}
\textbf{Componente} & \textbf{Puerto} & \textbf{Responsabilidad} \\
\hline
shell & 4300 & Punto de entrada, layout, autenticacion y navegacion. \\
mf-clientes & 4201 & Disponibilidad, creacion y consulta de reservas del cliente. \\
mf-administracion & 4202 & Gestion administrativa de canchas, usuarios y reservas. \\
mf-reportes & 4203 & Consulta de reportes basicos. \\
ms-usuarios & 8081 & Registro, login JWT, roles y usuarios. \\
ms-canchas & 8082 & Canchas, deportes, horarios y bloqueos. \\
ms-reservas & 8083 & Disponibilidad, reservas, cancelacion y reglas de negocio. \\
ms-reportes & 8084 & Agregacion y auditoria de reportes. \\
\end{longtable}

\section{Ejecucion con Docker Compose}

Desde la raiz del repositorio:

\begin{lstlisting}[style=terminal]
cp infra/.env.example infra/.env
docker compose -f infra/docker-compose.yml up --build
\end{lstlisting}

La primera ejecucion puede tardar varios minutos porque construye imagenes Java y Angular. Las siguientes ejecuciones reutilizan cache de Docker.

\section{URLs de acceso}

\begin{longtable}{p{5cm}p{8cm}}
\textbf{Servicio} & \textbf{URL} \\
\hline
Aplicacion web & \url{http://localhost:4300} \\
Swagger ms-usuarios & \url{http://localhost:8081/swagger-ui.html} \\
Swagger ms-canchas & \url{http://localhost:8082/swagger-ui.html} \\
Swagger ms-reservas & \url{http://localhost:8083/swagger-ui.html} \\
Swagger ms-reportes & \url{http://localhost:8084/swagger-ui.html} \\
\end{longtable}

\section{Credenciales de demostracion}

\begin{longtable}{p{4cm}p{5cm}p{4cm}}
\textbf{Rol} & \textbf{Correo} & \textbf{Clave} \\
\hline
Administrador & \texttt{admin@canchas.local} & \texttt{Admin123*} \\
\end{longtable}

Los usuarios cliente se crean desde la pantalla de registro. Las credenciales son solo para demostracion local; en otro ambiente deben reemplazarse las variables \texttt{BOOTSTRAP\_ADMIN\_PASSWORD}, \texttt{JWT\_SECRET} y \texttt{SERVICE\_API\_KEY}.

\section{Verificacion rapida}

\begin{enumerate}
  \item Abrir \url{http://localhost:4300}.
  \item Registrar un usuario cliente.
  \item Consultar disponibilidad de una cancha.
  \item Crear una reserva.
  \item Consultar mis reservas.
  \item Ingresar como administrador.
  \item Revisar canchas, usuarios, reservas y reportes.
  \item Abrir las cuatro interfaces Swagger.
\end{enumerate}

\evidenceimage{11-docker-compose.png}{Evidencia del entorno local desplegado con Docker Compose.}{fig:docker-compose}

\section{Apagado del entorno}

\begin{lstlisting}[style=terminal]
docker compose -f infra/docker-compose.yml down
\end{lstlisting}

Para eliminar tambien los volumenes de PostgreSQL y perder los datos locales:

\begin{lstlisting}[style=terminal]
docker compose -f infra/docker-compose.yml down -v
\end{lstlisting}

\section{Observaciones}

El sistema no utiliza API Gateway ni BFF. Los microfrontends consumen las API publicas de los microservicios y los microservicios se comunican entre si mediante endpoints internos protegidos con la cabecera \texttt{X-Service-Key}.

\end{document}

